% =============================================================================
%                         WEEK 2: IMAGE FORMATION
% =============================================================================
\section{Week 2: Image Formation}

\subsection{Topic Summary}

\begin{keybox}[Learning Objectives]
By the end of this week, you will be able to:
\begin{itemize}
  \item Explain the end-to-end process of digital image capture.
  \item Analyse the biological mechanisms of visual perception in the human eye.
  \item Compare the different strategies used by digital cameras and the human retina to efficiently encode spatial and colour information.
\end{itemize}
\end{keybox}

\textbf{Big Picture.}
This week covers how images are formed---both in cameras and in the eye. We examine the \term{physics} (light, reflectance, optics), \term{geometry} (projection, camera models), \term{digitisation} (sampling, quantization), and \term{biological vision} (photoreceptors, ganglion cells). The key insight is that recovering surface properties from images is ill-posed because $L = f(E, R)$---we observe luminance but need to infer reflectance without knowing illumination.

\separator

\textbf{What Examiners Like to Ask:}
\begin{itemize}
  \item Explain the relationship $L(\lambda) = f(E(\lambda), R(\lambda))$ and why colour constancy is hard.
  \item Derive the perspective projection equations and explain intrinsic/extrinsic parameters.
  \item Compare pinhole cameras vs lensed cameras (trade-offs).
  \item Explain demosaicing and different interpolation methods.
  \item Describe centre-surround receptive fields and their functional benefits.
  \item Compare foveal vs peripheral vision.
\end{itemize}

% =============================================================================
\subsection{Core Concepts}

\textbf{Roadmap:}
\begin{enumerate}
  \item Physics of image formation (light, reflectance, optics)
  \item Geometry of image formation (camera models, projective geometry)
  \item Digital images (digitisation, representation)
  \item Camera sensors (CCD, demosaicing)
  \item The eye (photoreceptors, ganglion cells)
\end{enumerate}

\bigseparator

% -----------------------------------------------------------------------------
\subsubsection{Physics of Image Formation}

\textbf{Overview of Image Formation:}

Images are formed when a \textbf{sensor} registers \textbf{radiation} that has interacted with \textbf{physical objects}. Light from a source is reflected by surfaces, passes through optics, and is recorded by a sensor.

\begin{center}
Light Source $\rightarrow$ Surface (Reflectance) $\rightarrow$ Optics $\rightarrow$ Sensor
\end{center}

\separator

\textbf{Two Sets of Parameters:}

\begin{center}
\renewcommand{\arraystretch}{1.3}
\begin{tabular}{@{} l p{10cm} @{}}
\toprule
\textbf{Parameter Type} & \textbf{What It Determines} \\
\midrule
Radiometric & Intensity/colour at each image location (illumination, surface reflectance, sensor sensitivity) \\
Geometric & Where scene points appear in image (camera position/orientation, optics, projection geometry) \\
\bottomrule
\end{tabular}
\end{center}

\separator

\textbf{Light and Colour:}

\begin{defbox}[Visible Light]
Electromagnetic radiation with wavelength $\lambda$ between 400--700nm. The sun's intensity peaks in this range, and human eyes evolved sensitivity to it.
\end{defbox}

\begin{itemize}
  \item \textbf{Illuminance} $E(x,y,\lambda)$: amount of light striking a surface (varies by light source)
  \item \textbf{Reflectance} $R(x,y,\lambda)$: fraction of light reflected at each wavelength (\term{albedo})
  \item \textbf{Luminance} $L(x,y,\lambda)$: amount of light reaching the sensor
\end{itemize}

\begin{thmbox}[Fundamental Equation of Image Formation]
\[
L(x,y,\lambda) = f\bigl(E(x,y,\lambda), R(x,y,\lambda)\bigr)
\]
The luminance depends on both illuminance and reflectance. This is \textbf{ill-posed}: we observe $L$ but need to recover $R$ without knowing $E$.
\end{thmbox}

\separator

\textbf{Colour Mixing:}

\begin{center}
\renewcommand{\arraystretch}{1.3}
\begin{tabular}{@{} l l l @{}}
\toprule
\textbf{Type} & \textbf{Process} & \textbf{Result} \\
\midrule
Additive (light) & R + G + B lights & White (broad spectrum) \\
Subtractive (pigment) & R + G + B pigments & Black (absorbs all) \\
\bottomrule
\end{tabular}
\end{center}

\begin{exbox}[Why Objects Have Colour]
An object looks \textbf{green} because it absorbs red and blue light, leaving more green in the reflected light. The colour we perceive depends on which wavelengths are \textit{not} absorbed.
\end{exbox}

\separator

\textbf{Colour Constancy:}

\begin{defbox}[Colour Constancy]
The ability to perceive surface colour as constant despite changes in illumination. Humans do this well; computers struggle.
\end{defbox}

\textbf{Approaches to estimate illumination $E$:}
\begin{itemize}
  \item Assume average reflectance across scene is known (grey world)
  \item Fix brightest patch to white
  \item Gamut falls within known range
  \item Known reference colour (colour chart, skin)
  \item Specular reflections have illumination colour
\end{itemize}

None work reliably in all cases. Human vision discounts illumination effects through learned priors.

\bigseparator

% -----------------------------------------------------------------------------
\subsubsection{Optics: Pinhole and Lens Cameras}

\begin{defbox}[Pinhole Camera]
A camera that restricts light flow through a small hole so that only one ray from each scene point reaches the sensor. Everything is in focus regardless of depth.
\end{defbox}

\textbf{Pinhole Trade-off:}
\begin{center}
\renewcommand{\arraystretch}{1.3}
\begin{tabular}{@{} l l l @{}}
\toprule
\textbf{Aperture} & \textbf{Focus} & \textbf{Brightness} \\
\midrule
Small pinhole & Sharp & Dim (long exposure, blur from motion) \\
Large pinhole & Blurred & Bright (short exposure) \\
\bottomrule
\end{tabular}
\end{center}

A \textbf{lens} solves this by focusing rays from the same point onto a single image location, keeping the image sharp while gathering more light.

\separator

\textbf{Thin Lens Properties:}
\begin{itemize}
  \item Rays through optical centre $O$ are not refracted
  \item Rays parallel to optical axis refract through focal point $F'$
  \item Rays through $F$ become parallel to optical axis
\end{itemize}

\begin{thmbox}[Thin Lens Equation]
\[
\frac{1}{f} = \frac{1}{|z|} + \frac{1}{|z'|}
\]
where $f$ = focal length, $z$ = object distance, $z'$ = image distance.
\end{thmbox}

\textbf{Key implications:}
\begin{itemize}
  \item Object at infinity $\Rightarrow$ image at $F'$
  \item Object at $F$ or closer $\Rightarrow$ no image formed
  \item \textbf{Depth of focus}: range of object depths that appear acceptably sharp
  \item Smaller aperture $\Rightarrow$ larger depth of focus (approaches pinhole behaviour)
\end{itemize}

\begin{tipbox}[Exam Tip: Lens vs Pinhole]
Know the trade-offs: pinhole has infinite depth of focus but needs long exposure; lens is bright but only objects at specific depths are in focus.
\end{tipbox}

\bigseparator

% -----------------------------------------------------------------------------
\subsubsection{Geometry of Image Formation}

\begin{defbox}[Perspective (Pinhole) Camera Model]
A mathematical model of geometric projection where a 3D point $P = (x,y,z)$ projects to image point $P' = (x',y',z')$. The image plane $\Pi'$ is at distance $f'$ from the pinhole along the optical axis.
\end{defbox}

\separator

\textbf{Equation of Projection:}

From similar triangles:
\[
\frac{x'}{x} = \frac{f'}{z} \quad\Rightarrow\quad x' = \frac{f'}{z} x
\]
\[
\frac{y'}{y} = \frac{f'}{z} \quad\Rightarrow\quad y' = \frac{f'}{z} y
\]

Since $P'$ lies on the image plane: $z' = f'$.

\textbf{In matrix form (homogeneous coordinates):}
\[
\begin{pmatrix} x' \\ y' \\ z' \end{pmatrix} = \frac{f'}{z} \begin{pmatrix} 1 & 0 & 0 & 0 \\ 0 & 1 & 0 & 0 \\ 0 & 0 & 1 & 0 \end{pmatrix} \begin{pmatrix} x \\ y \\ z \\ 1 \end{pmatrix}
\]

\separator

\textbf{Projective Geometry:}

\begin{center}
\renewcommand{\arraystretch}{1.3}
\begin{tabular}{@{} l p{5cm} p{5cm} @{}}
\toprule
\textbf{Geometry} & \textbf{Describes} & \textbf{Preserves} \\
\midrule
Euclidean & Objects ``as they are'' (3D$\rightarrow$3D) & Lengths, angles, parallelism \\
Projective & Objects ``as they appear'' (3D$\rightarrow$2D) & None of the above \\
\bottomrule
\end{tabular}
\end{center}

\textbf{Projective distortions:}
\begin{itemize}
  \item Distant objects appear smaller
  \item Parallel lines converge at \term{vanishing points}
  \item Equal lengths become unequal
  \item Right angles become non-right angles
\end{itemize}

\begin{defbox}[Vanishing Point]
The projection of a point at infinity. All parallel lines in 3D share the same vanishing point in the 2D image.
\end{defbox}

\begin{exbox}[Railway Tracks Example]
In the world: tracks parallel, ties equal length, perpendicular, evenly spaced.

In the image: tracks converge at vanishing point, ties get shorter with distance, not perpendicular to tracks, get closer together.
\end{exbox}

\bigseparator

% -----------------------------------------------------------------------------
\subsubsection{Intrinsic and Extrinsic Parameters}

\textbf{Converting Camera Coordinates to Pixels:}

\begin{defbox}[Intrinsic Parameters]
Camera-dependent properties that map the image plane [mm] to pixel coordinates [pixel]:
\begin{itemize}
  \item $o_x, o_y$: principal point (image centre) in pixels
  \item $\alpha = -f'/s_x$, $\beta = -f'/s_y$: magnification factors
  \item $s_x, s_y$: pixel size [mm/pixel]
  \item Aspect ratio = $\alpha/\beta$
\end{itemize}
\end{defbox}

\textbf{Pixel coordinates:}
\[
u = \alpha \frac{x}{z} + o_x, \qquad v = \beta \frac{y}{z} + o_y
\]

\textbf{In matrix form:}
\[
\begin{pmatrix} u \\ v \\ 1 \end{pmatrix} = \frac{1}{z} \begin{pmatrix} \alpha & 0 & o_x \\ 0 & \beta & o_y \\ 0 & 0 & 1 \end{pmatrix} \begin{pmatrix} 1 & 0 & 0 & 0 \\ 0 & 1 & 0 & 0 \\ 0 & 0 & 1 & 0 \end{pmatrix} \begin{pmatrix} x \\ y \\ z \\ 1 \end{pmatrix}
\]

\separator

\begin{defbox}[Extrinsic Parameters]
Camera location-dependent properties that transform world coordinates to camera coordinates:
\begin{itemize}
  \item $\mathbf{R}$: 3$\times$3 rotation matrix
  \item $\mathbf{t}$: 3$\times$1 translation vector
\end{itemize}
\end{defbox}

\textbf{World to camera transformation:}
\[
\begin{pmatrix} x \\ y \\ z \\ 1 \end{pmatrix} = \begin{pmatrix} \mathbf{R}^T & -\mathbf{R}^T\mathbf{t} \\ \mathbf{0}^T & 1 \end{pmatrix} \begin{pmatrix} X \\ Y \\ Z \\ 1 \end{pmatrix}
\]

\separator

\begin{thmbox}[Complete Projection Pipeline]
\[
\begin{pmatrix} u \\ v \\ 1 \end{pmatrix} = \frac{1}{z} \underbrace{\begin{pmatrix} \alpha & 0 & o_x \\ 0 & \beta & o_y \\ 0 & 0 & 1 \end{pmatrix}}_{\text{Intrinsic (2D$\rightarrow$2D)}} \underbrace{\begin{pmatrix} 1 & 0 & 0 & 0 \\ 0 & 1 & 0 & 0 \\ 0 & 0 & 1 & 0 \end{pmatrix}}_{\text{Projection (3D$\rightarrow$2D)}} \underbrace{\begin{pmatrix} \mathbf{R}^T & -\mathbf{R}^T\mathbf{t} \\ \mathbf{0}^T & 1 \end{pmatrix}}_{\text{Extrinsic (3D$\rightarrow$3D)}} \begin{pmatrix} X \\ Y \\ Z \\ 1 \end{pmatrix}
\]
\end{thmbox}

\begin{tipbox}[Exam Tip: Forward vs Inverse]
3D $\rightarrow$ 2D projection is \textbf{well-posed} (unique). 2D $\rightarrow$ 3D recovery is \textbf{ill-posed} (depth lost). Need extra information: stereo vision or scene priors.
\end{tipbox}

\bigseparator

% -----------------------------------------------------------------------------
\subsubsection{Digital Images}

\begin{defbox}[Digital Image]
A 2D array (matrix) of numbers $I(x,y)$ where each element represents the light intensity at a discrete location called a \term{pixel} (picture element).
\end{defbox}

\textbf{Two discretisation processes:}
\begin{itemize}
  \item \textbf{Pixelisation}: intensity values averaged within each grid location (spatial sampling)
  \item \textbf{Quantization}: intensity represented using finite discrete values (intensity sampling)
\end{itemize}

\separator

\textbf{Image Formats:}

\begin{center}
\renewcommand{\arraystretch}{1.3}
\begin{tabular}{@{} l l l l @{}}
\toprule
\textbf{Format} & \textbf{Values/Pixel} & \textbf{Bits} & \textbf{Range} \\
\midrule
Binary/Monochrome & 1 binary & 1 & $\{0, 1\}$ \\
Greyscale & 1 real & 8 & $[0, 255]$ or $[0, 1]$ \\
Colour (RGB) & 3 reals & 24 (8$\times$3) & 3 channels $\times$ 256 levels \\
\bottomrule
\end{tabular}
\end{center}

\textbf{Convention:} 0 = black, 255 (or 1.0) = white. Origin at top-left corner.

\separator

\textbf{Format Conversions:}

\textbf{RGB to Greyscale} (weighted average):
\[
I_{\text{grey}} = \frac{r \cdot I_R + g \cdot I_G + b \cdot I_B}{r + g + b}
\]

\textbf{Greyscale to Binary} (thresholding):
\[
I_{\text{bin}}(p) = \begin{cases} 1 & \text{if } I_{\text{grey}}(p) > t \\ 0 & \text{otherwise} \end{cases}
\]

\bigseparator

% -----------------------------------------------------------------------------
\subsubsection{Camera Sensors and Demosaicing}

\begin{defbox}[Charge Coupled Device (CCD)]
A semiconductor with a 2D matrix of photo-sensors. Each sensor accumulates charge proportional to incident light intensity and exposure time. Charge pattern corresponds to the image.
\end{defbox}

\textbf{Colour from CCDs:}

Photo-sensors are not wavelength-selective. Colour is achieved via filters:

\begin{center}
\renewcommand{\arraystretch}{1.3}
\begin{tabular}{@{} l p{4.5cm} p{4.5cm} @{}}
\toprule
\textbf{Solution} & \textbf{Pros} & \textbf{Cons} \\
\midrule
3 CCDs + prism & High quality, full sampling & Expensive, bulky \\
1 CCD + colour mask & Cheap, compact & Coarser sampling, needs demosaicing \\
\bottomrule
\end{tabular}
\end{center}

\separator

\begin{defbox}[Bayer Mask]
The most common colour filter array. Has 2$\times$ as many green filters as red/blue (matches human sensitivity). Results in subsampled colour channels that need interpolation.
\end{defbox}

\begin{defbox}[Demosaicing]
The process of filling in missing colour values so each pixel has R, G, and B values. A form of interpolation.
\end{defbox}

\textbf{Interpolation Methods:}

\begin{center}
\renewcommand{\arraystretch}{1.3}
\begin{tabular}{@{} l p{5cm} p{4cm} @{}}
\toprule
\textbf{Method} & \textbf{Description} & \textbf{Trade-off} \\
\midrule
Nearest Neighbour & Copy adjacent pixel from same channel & Fast, inaccurate \\
Bilinear & Average of 2 or 4 nearest same-channel pixels & Fast, blurs edges \\
Smooth Hue Transition & Interpolate ratio (hue) R/G or B/G & Better colour accuracy \\
Edge-Directed & Interpolate along axis with lowest gradient & Preserves edges \\
\bottomrule
\end{tabular}
\end{center}

\begin{exbox}[Edge-Directed Interpolation for Green]
To calculate $G_8$:
\begin{enumerate}
  \item Compute horizontal gradient: $\Delta_H = |G_7 - G_9|$
  \item Compute vertical gradient: $\Delta_V = |G_3 - G_{13}|$
  \item If $\Delta_H < \Delta_V$: $G_8 = (G_7 + G_9)/2$ (interpolate horizontally)
  \item If $\Delta_H > \Delta_V$: $G_8 = (G_3 + G_{13})/2$ (interpolate vertically)
  \item If equal: $G_8 = (G_3 + G_7 + G_9 + G_{13})/4$
\end{enumerate}
\end{exbox}

\bigseparator

% -----------------------------------------------------------------------------
\subsubsection{The Eye}

\textbf{Eye Anatomy:}
\begin{itemize}
  \item \textbf{Cornea}: bulk of refraction (fixed focus)
  \item \textbf{Lens}: further refraction, can change shape to adjust focal length
  \item \textbf{Iris}: regulates light amount (like camera aperture)
  \item \textbf{Retina}: contains photoreceptors
  \item \textbf{Optic nerve}: $\sim$1 million fibres transmit to brain
\end{itemize}

\separator

\textbf{Photoreceptor Types:}

\begin{center}
\renewcommand{\arraystretch}{1.3}
\begin{tabular}{@{} l l l l @{}}
\toprule
\textbf{Type} & \textbf{Sensitivity} & \textbf{Colour} & \textbf{Location} \\
\midrule
Rods & High (dim light) & None (monochrome) & Periphery \\
Cones (S) & Low (bright light) & Blue ($\sim$420nm) & Fovea \\
Cones (M) & Low (bright light) & Green ($\sim$540nm) & Fovea \\
Cones (L) & Low (bright light) & Red ($\sim$570nm) & Fovea \\
\bottomrule
\end{tabular}
\end{center}

\textbf{Distribution:}
\begin{itemize}
  \item \textbf{Fovea}: no rods, high cone density. \#blue $\ll$ \#red $\leq$ \#green
  \item \textbf{Periphery}: mostly rods, few cones
  \item \textbf{Blind spot}: no photoreceptors (optic nerve exit)
\end{itemize}

\separator

\textbf{Foveal vs Peripheral Vision:}

\begin{center}
\renewcommand{\arraystretch}{1.3}
\begin{tabular}{@{} l l l @{}}
\toprule
\textbf{Property} & \textbf{Fovea} & \textbf{Periphery} \\
\midrule
Resolution & High (dense receptors) & Low (sparse receptors) \\
Colour & Yes (cones) & No (rods) \\
Sensitivity & Low (cones) & High (rods) \\
Brain resources & Disproportionately large & Smaller \\
\bottomrule
\end{tabular}
\end{center}

\separator

\textbf{Retinal Processing---Ganglion Cells:}

\begin{defbox}[Receptive Field (RF)]
The area of visual space from which a neuron receives input. More generally, the stimulus properties that produce a response.
\end{defbox}

\begin{itemize}
  \item $\sim$1 million ganglion cells vs $\sim$100 million photoreceptors
  \item Each ganglion cell collects from multiple photoreceptors
  \item Large convergence in periphery, smaller in fovea (another reason for lower peripheral acuity)
\end{itemize}

\separator

\begin{defbox}[Centre-Surround Receptive Fields]
Ganglion cells have centre-surround RFs:
\begin{itemize}
  \item \textbf{On-centre, off-surround}: active if centre brighter than surround
  \item \textbf{Off-centre, on-surround}: active if centre darker than surround
\end{itemize}
\end{defbox}

\begin{thmbox}[Functional Benefits of Centre-Surround RFs]
\begin{enumerate}
  \item \textbf{Edge enhancement}: weak response on uniform regions, strong near edges
  \item \textbf{Invariance to ambient light}: measures contrast, not absolute intensity
  \item \textbf{Efficient coding}: only signals where intensity changes (``redundancy reduction'')
\end{enumerate}
\end{thmbox}

\begin{exbox}[Centre-Surround Response]
On a uniform surface: input to centre and surround cancels $\Rightarrow$ weak response.

At an edge: input unequal $\Rightarrow$ strong response (increased or decreased depending on type).

Under different lighting: if illumination doubles, both centre and surround double, so contrast unchanged $\Rightarrow$ response unchanged.
\end{exbox}

\separator

\textbf{Colour Opponent Cells:}

Ganglion cells combine cone inputs to create colour-opponent responses:
\begin{itemize}
  \item \textbf{Red-Green opponent}: Red ON / Green OFF (or vice versa)
  \item \textbf{Blue-Yellow opponent}: Blue ON / Yellow OFF (yellow = average of R and G cones)
\end{itemize}

\separator

\begin{defbox}[Difference of Gaussians (DoG)]
The standard mathematical model for centre-surround RFs:
\[
\text{DoG}(x,y) = G_{\sigma_1}(x,y) - G_{\sigma_2}(x,y)
\]
where $G_\sigma$ is a Gaussian with standard deviation $\sigma$, and $\sigma_1 < \sigma_2$.
\end{defbox}

\bigseparator

% =============================================================================
\subsection{Worked Examples}

\begin{exbox}[Example 1: Colour Constancy Problem]
\textbf{Q:} Why is recovering surface colour from an image an ill-posed problem?

\textbf{A:} The luminance $L(\lambda)$ we observe depends on both illuminance $E(\lambda)$ and reflectance $R(\lambda)$: $L = f(E, R)$. We want to recover $R$ (surface colour) but only measure $L$ and don't know $E$ (light source). Different combinations of $E$ and $R$ can produce the same $L$, making the problem ill-posed.
\end{exbox}

\begin{exbox}[Example 2: Thin Lens Calculation]
\textbf{Q:} A lens has focal length $f = 50$mm. An object is at distance $z = 200$mm. Where does the image form?

\textbf{A:} Using the thin lens equation:
\[
\frac{1}{f} = \frac{1}{|z|} + \frac{1}{|z'|} \quad\Rightarrow\quad \frac{1}{50} = \frac{1}{200} + \frac{1}{|z'|}
\]
\[
\frac{1}{|z'|} = \frac{1}{50} - \frac{1}{200} = \frac{4-1}{200} = \frac{3}{200}
\]
\[
|z'| = \frac{200}{3} \approx 66.7\text{mm}
\]
\end{exbox}

\begin{exbox}[Example 3: Perspective Projection]
\textbf{Q:} A point $P = (30, 40, 100)$ in camera coordinates. Focal length $f' = 10$mm. Find image coordinates.

\textbf{A:}
\[
x' = \frac{f'}{z} x = \frac{10}{100} \times 30 = 3\text{mm}
\]
\[
y' = \frac{f'}{z} y = \frac{10}{100} \times 40 = 4\text{mm}
\]
Image point: $P' = (3, 4, 10)$ mm.
\end{exbox}

\begin{exbox}[Example 4: Intrinsic Parameters]
\textbf{Q:} Convert camera coordinates $(x', y') = (3, 4)$mm to pixel coordinates. Given: $s_x = s_y = 0.01$mm/pixel, $f' = 10$mm, image centre $(o_x, o_y) = (320, 240)$.

\textbf{A:}
\[
\alpha = -\frac{f'}{s_x} = -\frac{10}{0.01} = -1000, \quad \beta = -\frac{f'}{s_y} = -1000
\]
From projection: $x'/z' = 3/10 = 0.3$, $y'/z' = 4/10 = 0.4$

Actually, using the formula directly with $x/z$:
\[
u = \alpha \frac{x}{z} + o_x = -1000 \times 0.3 + 320 = 20
\]
\[
v = \beta \frac{y}{z} + o_y = -1000 \times 0.4 + 240 = -160
\]
(Negative $v$ indicates point projects above image boundary.)
\end{exbox}

\begin{exbox}[Example 5: Centre-Surround Benefits]
\textbf{Q:} Explain why centre-surround RFs provide invariance to ambient lighting.

\textbf{A:} Centre-surround cells measure \textbf{contrast}---the difference between centre and surround intensities. If ambient light doubles, both centre and surround receive twice the light, but their \textit{difference} (relative to their sum) remains constant. The response depends on $I_{\text{centre}} - I_{\text{surround}}$, which encodes local contrast, not absolute intensity.
\end{exbox}

\bigseparator

% =============================================================================
\subsection{Summary}

\begin{keybox}[Key Takeaways]
\begin{enumerate}
  \item \textbf{Image formation}: $L = f(E, R)$---luminance depends on illumination and reflectance; recovering $R$ is ill-posed.
  \item \textbf{Optics}: Pinhole = sharp but dim; Lens = bright but limited depth of focus. Thin lens: $1/f = 1/|z| + 1/|z'|$.
  \item \textbf{Projection}: $x' = (f'/z)x$, $y' = (f'/z)y$. Parallel lines $\rightarrow$ vanishing points.
  \item \textbf{Camera parameters}: Intrinsic (focal length, pixel size, centre) + Extrinsic (rotation, translation).
  \item \textbf{Digital images}: Pixelisation (spatial) + Quantization (intensity). Bayer mask needs demosaicing.
  \item \textbf{Eye}: Fovea (high-res, colour, low-sensitivity) vs Periphery (low-res, mono, high-sensitivity).
  \item \textbf{Ganglion cells}: Centre-surround RFs provide edge enhancement, lighting invariance, efficient coding.
\end{enumerate}
\end{keybox}

\begin{center}
\renewcommand{\arraystretch}{1.3}
\begin{tabular}{@{} l l @{}}
\toprule
\textbf{Camera} & \textbf{Eye} \\
\midrule
Uniform pixel grid & Non-uniform (fovea vs periphery) \\
Bayer mask for colour & 3 cone types for colour \\
Demosaicing interpolates & Brain processes incomplete data \\
Raw pixel output & Centre-surround preprocessing \\
Fixed sensitivity & Rods (dim) vs Cones (bright) \\
\bottomrule
\end{tabular}
\end{center}
